% ID: SEL-D0021-OA-E04
% ESTADO: cotejada. Segunda pasada visual completa y borrador compilado revisados.
% Fuente: Selectividad/matematicas_aplicadas/reserva_a_EXAMEN 3-ANDALUCÍA.pdf
% SHA-256: 0b99df3e2c9a7af0a074136b72c3d257ad9241ccc79fed32894de424481b6be2
% Página: 1; opción A; ejercicio original 4.

Se desea estimar el porcentaje de alumnos de un determinado instituto que lleva gafas. Para ello se eligen 300 alumnos, de los que 210 llevan gafas.
\begin{enumerate}[label=\alph*)]
\item \textbf{(1.5 puntos)} Calcule el intervalo de confianza para la proporción de alumnos que lleva gafas, con un nivel de confianza del 97\%.
\item \textbf{(1 punto)} Si por estudios en otros institutos se sabe que la proporción de alumnos que lleva gafas es del 70\%, determine el tamaño mínimo de la muestra necesario para que, con una confianza del 97\%, el error máximo que se cometa sea inferior a $0.06$.
\end{enumerate}
